% Options for packages loaded elsewhere
\PassOptionsToPackage{unicode}{hyperref}
\PassOptionsToPackage{hyphens}{url}
\PassOptionsToPackage{dvipsnames,svgnames,x11names}{xcolor}
\documentclass[
]{article}
\usepackage{xcolor}
\usepackage[margin = 0.5in]{geometry}
\usepackage{amsmath,amssymb}
\setcounter{secnumdepth}{-\maxdimen} % remove section numbering
\usepackage{iftex}
\ifPDFTeX
  \usepackage[T1]{fontenc}
  \usepackage[utf8]{inputenc}
  \usepackage{textcomp} % provide euro and other symbols
\else % if luatex or xetex
  \usepackage{unicode-math} % this also loads fontspec
  \defaultfontfeatures{Scale=MatchLowercase}
  \defaultfontfeatures[\rmfamily]{Ligatures=TeX,Scale=1}
\fi
\usepackage{lmodern}
\ifPDFTeX\else
  % xetex/luatex font selection
\fi
% Use upquote if available, for straight quotes in verbatim environments
\IfFileExists{upquote.sty}{\usepackage{upquote}}{}
\IfFileExists{microtype.sty}{% use microtype if available
  \usepackage[]{microtype}
  \UseMicrotypeSet[protrusion]{basicmath} % disable protrusion for tt fonts
}{}
\makeatletter
\@ifundefined{KOMAClassName}{% if non-KOMA class
  \IfFileExists{parskip.sty}{%
    \usepackage{parskip}
  }{% else
    \setlength{\parindent}{0pt}
    \setlength{\parskip}{6pt plus 2pt minus 1pt}}
}{% if KOMA class
  \KOMAoptions{parskip=half}}
\makeatother
\usepackage{graphicx}
\makeatletter
\newsavebox\pandoc@box
\newcommand*\pandocbounded[1]{% scales image to fit in text height/width
  \sbox\pandoc@box{#1}%
  \Gscale@div\@tempa{\textheight}{\dimexpr\ht\pandoc@box+\dp\pandoc@box\relax}%
  \Gscale@div\@tempb{\linewidth}{\wd\pandoc@box}%
  \ifdim\@tempb\p@<\@tempa\p@\let\@tempa\@tempb\fi% select the smaller of both
  \ifdim\@tempa\p@<\p@\scalebox{\@tempa}{\usebox\pandoc@box}%
  \else\usebox{\pandoc@box}%
  \fi%
}
% Set default figure placement to htbp
\def\fps@figure{htbp}
\makeatother
\setlength{\emergencystretch}{3em} % prevent overfull lines
\providecommand{\tightlist}{%
  \setlength{\itemsep}{0pt}\setlength{\parskip}{0pt}}
\usepackage{color}
\usepackage{framed}
\setlength{\fboxsep}{.8em}

\usepackage{tcolorbox}

\tcbset{
  mystyle/.style={
    before skip=-10pt, % Adjust this value
    after skip=-10pt   % Adjust this value
  }
}
\newtcolorbox{mynoteblock}[1][]{mystyle, #1}

\newtcolorbox{blackbox}{
  colback=black,
  colframe=orange,
  coltext=white,
  boxsep=5pt,
  arc=4pt}
  
\newtcolorbox{ltbluebox}{
  colback=blue!5!white,
  colframe=blue!75!black,
  boxsep=5pt,
  arc=4pt}
  
\newtcolorbox{ltgreenbox}{
  colback=green!5!white,
  colframe=blue!75!black,
  boxsep=5pt,
  arc=4pt}
  
\newtcolorbox{ltpurplebox}{
  colback=purple!5!white,
  colframe=blue!75!black,
  boxsep=5pt,
  arc=4pt}
  


\newenvironment{cols}[1][]{}{}

\newenvironment{col}[1]{\begin{minipage}[t]{#1}\ignorespaces\vspace{0pt}}{%
\end{minipage}
\ifhmode\unskip\fi
\aftergroup\useignorespacesandallpars}

\def\useignorespacesandallpars#1\ignorespaces\fi{%
#1\fi\ignorespacesandallpars}

\makeatletter
\def\ignorespacesandallpars{%
  \@ifnextchar\par
    {\expandafter\ignorespacesandallpars\@gobble}%
    {}%
}
\makeatother

\usepackage{awesomebox}
\setlength{\aweboxvskip}{1mm}


\pagenumbering{gobble}

\usepackage{enumitem}
\setlist{itemsep=1pt, parsep=0pt} %# Adjust vertical spacing
\setlist[itemize,1]{leftmargin=*, labelsep=0.3em} %# Adjust horizontal spacing

\usepackage{bookmark}
\IfFileExists{xurl.sty}{\usepackage{xurl}}{} % add URL line breaks if available
\urlstyle{same}
\hypersetup{
  colorlinks=true,
  linkcolor={Maroon},
  filecolor={Maroon},
  citecolor={Blue},
  urlcolor={blue},
  pdfcreator={LaTeX via pandoc}}

\author{}
\date{\vspace{-2.5em}}

\begin{document}

\vspace{-1cm}

\section{MAFMC Indicators Project:
Definitions}\label{mafmc-indicators-project-definitions}

\subsection{What is an ecosystem
indicator?}\label{what-is-an-ecosystem-indicator}

\begin{noteblock}
An \emph{indicator} tells us something about where we are relative to
our goals or to limits, or about the context we are working within that
may affect achieving our goals.

\end{noteblock}

\includegraphics[width=1\linewidth]{WK1_Definitions_files/figure-latex/soefigs-1}

\subsection{What is a SMART indicator?}\label{what-is-a-smart-indicator}

\begin{noteblock}

\begin{itemize}
\tightlist
\item
  \textbf{Specific}: clearly defined, describing exactly what is
  measured at what scale and how the indicator is calculated so that
  methods can be repeated. The indicator can be directly linked to a
  management process objective.
\end{itemize}

\end{noteblock}

\begin{noteblock}

\begin{itemize}
\tightlist
\item
  \textbf{Measurable}: appropriate quantitative or qualitative data are
  available to evaluate change over time or in space. Tied to measurable
  management process objective.
\end{itemize}

\end{noteblock}

\begin{noteblock}

\begin{itemize}
\tightlist
\item
  \textbf{Achievable}: responses to changes in management are detectable
  within the management time frame--objectives tied to the indicator
  must be achievable.
\end{itemize}

\end{noteblock}

\begin{noteblock}

\begin{itemize}
\tightlist
\item
  \textbf{Relevant}: accurately reflects changes in the underlying
  process of concern. The indicator can be directly linked to a
  management process objective.
\end{itemize}

\end{noteblock}

\begin{noteblock}

\begin{itemize}
\tightlist
\item
  \textbf{Time-bound}: how often is indicator measured and does this
  align with the process of concern, and the management process?
\end{itemize}

\end{noteblock}

SMART indicators are Specific, Measurable, Achievable, Relevant, and
Timely for the priority Council decisions. A SMART indicator for one
process won't necessarily be SMART for another.

\begin{ltgreenbox}

\subsection{\texorpdfstring{What is Operational? \emph{In use or ready
for
use}}{What is Operational? In use or ready for use}}\label{what-is-operational-in-use-or-ready-for-use}

Components of an operational process:

\begin{itemize}
\item
  Clear articulation of management objectives (use)
\item
  Standardized information updates on a regular reporting schedule
\item
  Feedbacks from the management system and adjustments to achieve
  management objectives
\end{itemize}

``Operationalizing'' is an iterative process within a specific context

\end{ltgreenbox}

\end{document}
